\documentclass[12pt,a4paper]{article}
\usepackage[margin=25mm, headheight=15pt, headsep=10pt]{geometry}
\usepackage{fontspec}
\usepackage[table]{xcolor} % Note: table option is required for rowcolors
\usepackage{titlesec}
\usepackage{tocloft}
\usepackage{fancyhdr}
\usepackage{booktabs}
\usepackage{tcolorbox}
\tcbuselibrary{skins, breakable}
\usepackage[colorlinks=true, linkcolor=emerald, urlcolor=emerald, bookmarks=true]{hyperref}

\definecolor{emerald}{RGB}{0,102,51}
\definecolor{darkred}{RGB}{139,0,0}
\definecolor{lightgray}{RGB}{245,245,245}

\usepackage[nil]{babel}
\babelprovide[import, main]{bengali}
\babelprovide[import]{arabic}
\babelfont{rm}[Renderer=HarfBuzz,Scale=1.0]{Noto Serif Bengali}
\babelfont{sf}[Renderer=HarfBuzz]{Noto Sans Bengali}
\babelfont{tt}[Renderer=HarfBuzz]{Noto Sans Bengali}
\babelfont[arabic]{rm}[Renderer=HarfBuzz,Scale=1.1]{Noto Naskh Arabic}

% Section styling
\titleformat{\section}{\Large\bfseries\color{emerald}}{\thesection.}{0.5em}{}
\titleformat{\subsection}{\large\bfseries\color{emerald!85!black}}{\thesubsection.}{0.5em}{}
\titleformat{\subsubsection}{\normalsize\bfseries\color{emerald!70!black}}{\thesubsubsection.}{0.5em}{}
\titlespacing*{\section}{0pt}{18pt}{8pt}

% Fancy Header/Footer setup
\pagestyle{fancy}
\fancyhf{} % clear all header and footer fields
\fancyhead[L]{\color{emerald}\small\textbf{\nouppercase{\leftmark}}}
\fancyfoot[L]{\scriptsize\color{black!50}{\lat{AI}}-সংকলিত, আলেম-যাচাইকৃত নয়}
\fancyfoot[C]{\color{emerald!70!black}\thepage}
\renewcommand{\headrulewidth}{0.4pt}
\renewcommand{\headrule}{\hbox to\headwidth{\color{emerald}\leaders\hrule height \headrulewidth\hfill}}
\renewcommand{\footrulewidth}{0pt}

% Custom boxes
% 1. Ayat/Hadith Box
\newtcolorbox{ayatbox}[1][]{
    enhanced, breakable, 
    colback=emerald!5, 
    colframe=emerald, 
    boxrule=0.5pt, 
    arc=3pt, 
    left=10pt, right=10pt, top=10pt, bottom=10pt,
    #1
}

% 2. Quote/Translation Box
\newtcolorbox{quotebox}[1][]{
    enhanced, breakable,
    frame hidden,
    colback=lightgray,
    borderline west={3pt}{0pt}{emerald},
    left=15pt, right=10pt, top=10pt, bottom=10pt,
    sharp corners,
    #1
}

% 3. AI-generated notice box
\newtcolorbox{warnbox}[1][]{
    enhanced, breakable,
    colback=red!4,
    colframe=darkred,
    boxrule=0.8pt,
    arc=3pt,
    left=10pt, right=10pt, top=10pt, bottom=10pt,
    #1
}

% Commands
\newcommand{\arab}[1]{\foreignlanguage{arabic}{#1}}
\newcommand{\kib}[1]{\textcolor{emerald}{\textbf{#1}}}

% Latin (English) fallback: Noto Serif Bengali has NO Latin glyphs
\newfontfamily\latfont[Renderer=HarfBuzz]{Noto Serif}
\newcommand{\lat}[1]{{\latfont #1}}

\setlength{\parskip}{5pt}
\setlength{\parindent}{0pt}

% Fix for missing bullet in Noto Serif Bengali
\renewcommand{\labelitemi}{$\bullet$}



\begin{document}
\begin{center}{\Large\bfseries\color{emerald} ক্ষমা করা পাপ — কিয়ামতের দিন কী হয়? (ব্যবহারিক গাইড)}\end{center}
\vspace{1em}
\section*{ক্ষমা করা পাপ — কিয়ামতের দিন কী হয়? (ব্যবহারিক গাইড)}

\begin{warnbox}
 \textbf{গুরুত্বপূর্ণ নোটিশ (\lat{Important} \lat{Notice}):} এই কন্টেন্টটি \lat{AI} (কৃত্রিম বুদ্ধিমত্তা) দিয়ে প্রস্তুত ও সংকলিত। \lat{AI} কঠোর সূত্র-মান নীতি মেনে শুধু নির্ভরযোগ্য উৎস থেকে তথ্য সংগ্রহ করেছে — কেবল সহীহ/হাসান হাদিস ও সহীহ সনদের আসার রাখা হয়েছে, দুর্বল সনদ বাদ দেওয়া হয়েছে। তবে এটি কোনো আলেম বা মুহাদ্দিস কর্তৃক যাচাইকৃত নয়।
\end{warnbox}

\medskip\hrule\medskip

\subsection*{প্রশ্নটির প্রকৃত রূপ}

তাওবাকারী বা পাপী মানুষের মনে যে ভয়টি আসে সেটি এমন:

\begin{quote}
\lat{“}আল্লাহ যদি আমার এই পাপ ক্ষমা করে দেন — তাহলে কিয়ামতের দিন কি সেটা আবার তোলা হবে? আমলনামায় থাকবে কি? সবার সামনে প্রকাশ হবে কি? আমার কি আবার জিজ্ঞাসাবাদ হবে?\lat{”}
\end{quote}

এই ভয়ের মূলে আছে তিনটি জিনিসের মধ্যে গুলিয়ে ফেলা: \textbf{ক্ষমা (মাগফিরাহ), গোপন রাখা (সতর) ও হিসাব (হিসাব)}। এগুলো একই জিনিস নয় — এবং কুরআন-হাদিসে প্রতিটির নিজস্ব বিধান আছে। সেগুলো আলাদা করে বোঝাই এই দ্বিধা দূর করার চাবিকাঠি।

\medskip\hrule\medskip

\subsection*{১. \lat{“}ক্ষমা\lat{”} মানে আসলে কী?}

\subsubsection*{শব্দের ভিত্তি}

\begin{itemize}
\item \textbf{মাগফিরাহ (\arab{مغفرة})} — মূল ধাতু \lat{“}গাফর\lat{”} — অর্থ ঢেকে দেওয়া, আবৃত করা, রক্ষা করা। ক্ষমা মানে পাপের শাস্তি ও প্রভাব \textbf{ঢেকে দেওয়া}।
\item \textbf{সতর (\arab{ستر})} — গোপন রাখা, পর্দা করা।
\item \textbf{মাহু (\arab{محو})} — মুছে ফেলা।
\end{itemize}
\subsubsection*{দলিলের আলোকে ক্ষমার তিনটি স্তর}

\begin{table}[h]\centering\small
\begin{tabular}{lll}
স্তর & অর্থ & দলিল \\
\hline
\textbf{সতর (গোপন)} & আল্লাহ পাপ ঢেকে রাখেন — দুনিয়ায় ও আখিরাতে & বুখারি ২৪৪১ — \lat{“}আমি দুনিয়ায় এগুলো তোমার উপর গোপন রেখেছিলাম\lat{”}; মুসলিম ২৬৯৯ \\
\textbf{মাহু (মুছে দেওয়া)} & তাওবার পর পাপের রেকর্ড মুছে যায় / নেকিতে বদলায় & ২৫:৭০ — \lat{“}পাপগুলো নেকিতে বদলে দেব\lat{”}; ১১:১১৪ — \lat{“}ভালো কাজ মন্দ কাজ মুছে দেয়\lat{”} \\
\textbf{মাগফিরাহ (ক্ষমা)} & শাস্তি না হওয়া — হিসাবের বোঝা না থাকা & বুখারি ৬৫৩৬ — সহজ হিসাব = আরয; বুখারি ২৪৪১ — ক্ষমার পর নেকির আমলনামা \\
\hline
\end{tabular}
\end{table}

\textbf{মূল কথা:} তাওবা করলে পাপ শুধু \lat{“}মাফ\lat{”} হয় না — আল্লাহ তা \textbf{মুছে দেন বা নেকিতে বদলে দেন} (২৫:৭০), এবং \textbf{গোপনও রাখেন} (বুখারি ২৪৪১)। তিনটি কাজই আল্লাহর পক্ষ থেকে।

\medskip\hrule\medskip

\subsection*{২. ক্ষমা করা পাপ কি আমলনামায় থাকবে?}

\subsubsection*{সরাসরি উত্তর}

\textbf{তাওবার পর পাপ আর আমলনামার হিসাবের অংশ থাকে না।} দলিল:

\begin{enumerate}
\item \textbf{২৫:৭০:} \lat{“}তাদের পাপগুলো আল্লাহ নেকিতে বদলে দেবেন\lat{”} — পাপের জায়গায় নেকি লেখা হয়। ইবনে আব্বাস (রাঃ)-এর ব্যাখ্যা অনুযায়ী আক্ষরিক অর্থেই রেকর্ড পরিবর্তিত হয়; অপর ব্যাখ্যায় শাস্তির বদলে ক্ষমা ও নেকি — উভয় ক্ষেত্রেই পাপ আর বোঝা নয়।
\item \textbf{নাজওয়া হাদিস (বুখারি ২৪৪১):} ক্ষমার পর মুমিনকে \textbf{\lat{“}নেকির আমলনামা\lat{”}} দেওয়া হয় — পাপের আমলনামা নয়।
\item \textbf{১১:১১৪:} \lat{“}ভালো কাজ মন্দ কাজ মুছে দেয়\lat{”} — নেক আমল পাপকে ঢেকে/মুছে দেয়।
\end{enumerate}
\subsubsection*{যেটুকু সতর্কতা}

\begin{itemize}
\item ১৮:৪৯-তে \lat{“}আমলনামা ছোট-বড় কিছুই বাদ দেয় না\lat{”} — এটি \textbf{তাওবাহীন অপরাধীর} বর্ণনা (আয়াতের শব্দ: আল-মুজরিমীন)। তাওবাকারী ২৫:৭০-এর আওতায় — পাপ নেকিতে বদল।
\item আয়েশা (রাঃ)-এর প্রশ্নের জবাবে নবীজি (সা.) বললেন: সহজ হিসাব = \textbf{আরয (পেশকরণ)} — অর্থ মুমিনের আমল পেশ হবে, কিন্তু পুঙ্খানুপুঙ্খ জিজ্ঞাসাবাদ হবে না (বুখারি ৬৫৩৬)।
\end{itemize}
\medskip\hrule\medskip

\subsection*{৩. আবার জিজ্ঞাসাবাদ হবে কি?}

\subsubsection*{দুই ধরনের \lat{“}হিসাব\lat{”} আছে}

\begin{table}[h]\centering\small
\begin{tabular}{llll}
ধরন & কী হয় & কার জন্য & দলিল \\
\hline
\textbf{আরয (পেশকরণ)} & আমল দেখানো হয় — হিসাব সহজ & মুমিন — যার পাপ ক্ষমা হয়ে গেছে & বুখারি ৬৫৩৬ — \lat{“}সেটি তো আরয\lat{”} \\
\textbf{মুনাকাশাহ (পুঙ্খানুপুঙ্খ জিজ্ঞাসাবাদ)} & প্রতিটি কাজের খুঁটিনাটি তোলা & তাওবাহীন — শাস্তির কারণ & বুখারি ৬৫৩৬ — \lat{“}যার হিসাব পুঙ্খানুপুঙ্খ নেওয়া হবে সে শাস্তি পাবে\lat{”} \\
\hline
\end{tabular}
\end{table}

\subsubsection*{নাজওয়া হাদিসে \lat{“}পাপ স্মরণ করানো\lat{”} কীভাবে বুঝব?}

নাজওয়া হাদিসে (বুখারি ২৪৪১) আল্লাহ মুমিনকে কাছে টেনে তাঁর \textbf{পাপগুলো একান্তে স্মরণ করিয়ে} স্বীকার করান — তারপর ক্ষমা করেন। এটি কি \lat{“}আবার জিজ্ঞাসাবাদ\lat{”}?

\begin{itemize}
\item \textbf{না — এটি জিজ্ঞাসাবাদ নয়, এটি দয়ার নাজওয়া (গোপন আলাপ)।} এখানে:
\item পাপ \textbf{একান্তে} — কানাফ (আড়াল) দিয়ে — কোনো ফেরেশতা বা সৃষ্টিও দেখে না;
\item উদ্দেশ্য \textbf{শাস্তি নয়} — বান্দা নিজের অবস্থা বুঝুক ও আল্লাহর অনুগ্রহের মাত্রা অনুভব করুক;
\item পরিণতি \textbf{ক্ষমা} — \lat{“}আজ আমি তোমার জন্য ক্ষমা করে দিলাম\lat{”} — তারপর \textbf{নেকির আমলনামা}।
\item ইবনে হাজার (রহ.) বলেন: এই হাদিস প্রমাণ করে আল্লাহ মুমিনদের পাপ আখিরাতে \textbf{গোপন রাখবেন} — কাফির-মুনাফিকদের ক্ষেত্রে এই গোপন রাখা নেই, তাদের প্রকাশ করা হবে।
\item অর্থাৎ পাপের \lat{“}মেনশন\lat{”} হবে — কিন্তু \textbf{শাস্তি বা অপমানের জন্য নয়; বরং ক্ষমার ঘোষণার জন্য।} বান্দা যা ভেবেছিল ধ্বংস হবে — সেটাই পরিণত হয় নাজাতের সুসংবাদে।
\end{itemize}
\medskip\hrule\medskip

\subsection*{৪. সবার সামনে প্রকাশ হবে কি?}

\subsubsection*{কাদের পাপ প্রকাশ হবে}

\begin{table}[h]\centering\small
\begin{tabular}{lll}
ব্যক্তি & প্রকাশের ধরন & দলিল \\
\hline
\textbf{কাফির-মুনাফিক} & সবার সামনে ঘোষণা — \lat{“}এরা রবের বিরুদ্ধে মিথ্যা বলেছিল\lat{”} & বুখারি ২৪৪১; মুসলিম ২৭৬৮ \\
\textbf{বিশ্বাসঘাতক (গাদির)} & পতাকা উঁচু করা হয় — বিশ্বাসঘাতকতা চেনা যায় & বুখারি ৩১৮৮ \\
\textbf{যে মানুষের হক জুলুম করেছে} & নেকি থেকে হক আদায় — হিসাবের মাধ্যমে প্রকাশ & মুসলিম ২৫৮১; বুখারি ২৪৪৯ \\
\hline
\end{tabular}
\end{table}

\subsubsection*{কাদের পাপ গোপন থাকবে}

\begin{itemize}
\item \textbf{মুমিন তাওবাকারী} — নাজওয়া হাদিস: কানাফ (আড়াল) দিয়ে একান্তে ক্ষমা; প্রকাশ হয় না (বুখারি ২৪৪১)।
\item \textbf{যে দুনিয়ায় পাপ গোপন রেখেছে} — \lat{“}যে মুসলিমের দোষ গোপন রাখে, আল্লাহ দুনিয়া ও আখিরাতে তার দোষ গোপন রাখবেন\lat{”} (মুসলিম ২৬৯৯); \lat{“}সারারাত তার রব তাকে ঢেকে রেখেছিলেন\lat{”} (বুখারি ৬০৬৯)।
\end{itemize}
\subsubsection*{৮৬:৯-এর \lat{“}গোপন প্রকাশ\lat{”} কীভাবে বুঝব?}

\lat{“}যেদিন গোপন বিষয়সমূহ পরীক্ষা করা হবে\lat{”} (৮৬:৯) — অর্থ আল্লাহর কাছে সব প্রকাশ, কোনো কিছু লুকানো নেই। কিন্তু এটি \textbf{সবার সামনে প্রদর্শনের} প্রতিশ্রুতি নয় — নাজওয়া হাদিসে মুমিনের ক্ষেত্রে গোপন রাখার কথা স্পষ্ট। দুই আয়াত-হাদিসের সমন্বয়: \textbf{আল্লাহ সব জানেন ও হিসাব নেবেন — কিন্তু মুমিনের ক্ষমা করা পাপকে তিনি নিজেই ঢেকে রাখবেন।}

\medskip\hrule\medskip

\subsection*{৫. আল্লাহর হক বনাম মানুষের হক — মূল পার্থক্য}

এটি সবচেয়ে গুরুত্বপূর্ণ পার্থক্য, যা প্রশ্নটির অর্ধেক সমাধান করে:

\begin{table}[h]\centering\small
\begin{tabular}{lll}
বিষয় & আল্লাহর হক & মানুষের হক \\
\hline
\textbf{কীভাবে ক্ষমা হয়} & তাওবা, ইস্তিগফার, নেক আমল — ২৫:৭০; ১১:১১৪ & ফেরত দেওয়া/ক্ষমা আদায় করা ছাড়া নয় — মুসলিম ২৫৮১ \\
\textbf{আল্লাহর ক্ষমা} & সরাসরি — আল্লাহ ক্ষমা করেন & আল্লাহর ক্ষমা যথেষ্ট নয় — বান্দার ক্ষমা/ফেরত লাগবে \\
\textbf{কিয়ামতের হিসাব} & ক্ষমা হয়ে গেলে আর তোলা হয় না & হক আদায় না হলে কিয়ামতে নেকি থেকে কেটে নেওয়া হয় — মুসলিম ২৫৮১; বুখারি ২৪৪৯ \\
\textbf{উদাহরণ} & নামাজ ছুটে যাওয়া, গুনাহ করা & গালি দেওয়া, সম্পদ খাওয়া, অপবাদ, মারধর \\
\hline
\end{tabular}
\end{table}

\textbf{ব্যবহারিক নিয়ম:} তাওবার সময় নিজেকে প্রশ্ন করুন — \lat{“}এই পাপে কি কারো হক নষ্ট হয়েছে?\lat{”} যদি হ্যাঁ, তাহলে তা ফেরত দিন বা ক্ষমা আদায় করুন। আল্লাহর হকের ক্ষেত্রে তাওবা-ইস্তিগফারই যথেষ্ট।

\medskip\hrule\medskip

\subsection*{৬. কোন পাপের ক্ষেত্রে কী প্রযোজ্য — সারসংক্ষেপ}

\begin{table}[h]\centering\small
\begin{tabular}{ll}
পরিস্থিতি & কিয়ামতের দিন কী হবে \\
\hline
\textbf{তাওবা করা হয়েছে — আল্লাহর হক} & পাপ মুছে যায়/নেকিতে বদলায় — আরয হয়, জিজ্ঞাসাবাদ নয়; গোপন থাকে (২৫:৭০; বুখারি ২৪৪১; ৬৫৩৬) \\
\textbf{তাওবা করা হয়েছে — মানুষের হক ফেরত/ক্ষমা আদায় হয়েছে} & একই — ক্ষমা সম্পূর্ণ; হিসাব সহজ \\
\textbf{তাওবা করা হয়েছে — মানুষের হক ফেরত হয়নি} & হক আদায়ের হিসাব হবে — নেকি থেকে কেটে নেওয়া হতে পারে (মুসলিম ২৫৮১) \\
\textbf{তাওবা হয়নি — শিরকে মৃত্যু} & ক্ষমার বাইরে (৪:৪৮) \\
\textbf{তাওবা হয়নি — শিরক ছাড়া অন্য পাপ} & আল্লাহর ইচ্ছাধীন — তিনি চাইলে ক্ষমা (৪:৪৮); হিসাব হবে \\
\textbf{নতুন মুসলিম} & ইসলাম পূর্বের সব পাপ মুছে দেয় (মুসলিম ১২১) \\
\textbf{ভুল-ভুলে যাওয়া-বাধ্য} & মূলত ক্ষমা — হিসাবে তোলা হয় না (ইবনে মাজাহ ২০৪৫) \\
\textbf{পাপ প্রকাশ করে বেড়ানো (মুজাহির)} & গোপন রাখার সুযোগ হারায় (বুখারি ৬০৬৯) \\
\hline
\end{tabular}
\end{table}

\medskip\hrule\medskip

\subsection*{৭. পাপী ও তাওবাকারীর জন্য আশা — কীভাবে অযথা ভয় না পেয়ে আল্লাহর দিকে ফিরবেন}

\subsubsection*{ক. অতীতের ভয় যেন তাওবায় বাধা না হয়}

\begin{itemize}
\item আল্লাহর রহমত থেকে নিরাশ হওয়া নিষিদ্ধ — \lat{“}তোমরা আল্লাহর রহমত থেকে নিরাশ হয়ো না; আল্লাহ সব গুনাহ ক্ষমা করেন\lat{”} (৩৯:৫৩)।
\item তাওবার দরজা খোলা — প্রাণ কণ্ঠাগত হওয়ার আগ পর্যন্ত (তিরমিযি ৩৫৩৭ — হাসান)।
\item নাজওয়া হাদিসের বার্তা: \textbf{আল্লাহ কিয়ামতের দিনও পাপ গোপন রাখেন — দুনিয়ায় গোপন রাখার বিনিময়ে।} তাই পাপ গোপন রেখে, অনুতপ্ত হয়ে আল্লাহর দিকে ফেরা — কিয়ামতের ভয়কে আশায় বদলে দেয়।
\end{itemize}
\subsubsection*{খ. তাওবার পর করণীয়}

\begin{enumerate}
\item \textbf{পাপ ছেড়ে দিন ও আফসোস করুন} — \lat{“}অনুতাপই তাওবা\lat{”} (ইবনে মাজাহ ৪২৫২ — হাসান)।
\item \textbf{নেক আমল দিয়ে পাপ ঢেকে দিন} — \lat{“}ভালো কাজ মন্দ কাজ মুছে দেয়\lat{”} (১১:১১৪); নবীজি (সা.): \lat{“}পাপের পর ভালো কাজ করো — তা মুছে দেবে\lat{”} (তিরমিযি ১৯৮৭ — হাসান-সহীহ)।
\item \textbf{মানুষের হক ফেরত দিন} — মুসলিম ২৫৮১; বুখারি ২৪৪৯।
\item \textbf{পাপ প্রকাশ করবেন না} — গোপন রাখা আখিরাতের গোপন রাখার সাথে যুক্ত (বুখারি ৬০৬৯; মুসলিম ২৬৯৯)।
\item \textbf{আল্লাহর কাছে দোয়া করুন} — \lat{“}হে আল্লাহ! তুমি ক্ষমাশীল, ক্ষমা ভালোবাসো — আমাকে ক্ষমা করো\lat{”} (তিরমিযি ৩৫১৩ — সহীহ)।
\end{enumerate}
\subsubsection*{গ. চূড়ান্ত আশ্বাস}

কুরআন-হাদিসের সামগ্রিক চিত্র এমন:

\begin{quote}
\textbf{তাওবা করলে আল্লাহ পাপ মুছে দেন, নেকিতে বদলে দেন, গোপন রাখেন — কিয়ামতের দিন তাওবাকারীকে তার নেকির আমলনামা দেওয়া হয়; তার হিসাব সহজ (আরয); তার পাপ একান্তে স্মরণ করিয়ে ক্ষমা করা হয় — সবার সামনে নয়। শুধু শিরকে মৃত্যু ও অপরের হক নষ্ট — এই দুটির ক্ষেত্রে ভিন্ন বিধান; বাকি সব পাপের দরজা তাওবার মাধ্যমে খোলা।}
\end{quote}

অতএব — অতীতের পাপের ভয়ে হতাশ হওয়া নয়; বরং এখনই তাওবা, নেক আমল ও হক আদায়ের মাধ্যমে আল্লাহর দিকে ফিরে আসা — এটি কিয়ামতের ভয়কে শান্তি ও আশায় বদলে দেওয়ার একমাত্র ইসলামি পথ।

\medskip\hrule\medskip

\subsection*{মূল স্মরণীয় পয়েন্ট}

\begin{enumerate}
\item \textbf{ক্ষমা = ঢেকে দেওয়া + মুছে দেওয়া + নেকিতে বদল} — ২৫:৭০; বুখারি ২৪৪১।
\item \textbf{সহজ হিসাব = আরয (পেশকরণ)} — পুঙ্খানুপুঙ্খ জিজ্ঞাসাবাদ নয় — বুখারি ৬৫৩৬।
\item \textbf{মুমিনের পাপ গোপন থাকে} — প্রকাশ হয় কাফির-মুনাফিক ও বিশ্বাসঘাতকের — বুখারি ২৪৪১; ৩১৮৮।
\item \textbf{মানুষের হক ফেরত ছাড়া ক্ষমা নয়} — মুসলিম ২৫৮১।
\item \textbf{অতীতের ভয় তাওবায় বাধা নয়} — ৩৯:৫৩; তিরমিযি ৩৫৩৭।
\end{enumerate}

\end{document}
